%% kapa: exam=1 round=19 year=2008 period=2 form=A subject=english subject_ko=영어 seq=33 number=73
%% source: Q-Net 감정평가사 기출문제 (cst00309) artlSeq=1001962 file=extracted/1008/1001962/2008년 제19회 감정평가사 1차시험 기출문제/2교시.hwp
%% conversion: kordoc 4.24.51 -> tools/gen1.py (auto); review=proofread-sonnet; flags=none
\begin{question}{73}
다음 글을 읽고 밑줄 친 \uline{This book}에 대한 설명으로 가장 알맞은 것을 고르시오.
\begin{framedtext}
\uline{This book}, published in 1929, is oblique rather than polemical in its arguments. Based on speeches written for delivery at the British women's colleges of Girton and Newnham, it surpasses any feminist writing since Mary Wollstonecraft for artistry and profundity of emotion. Through her flexible and evocative use of language, the author conveys the longing, frustrating and disabling rages that are experienced by every woman and by the female creative writer especially. She looks ahead to a time when Shakespeare's sister might "find it possible to live and write her poetry."
\end{framedtext}
\begin{choices}
  \item a collection of political speeches
  \item an inquiry into language use
  \item a writing on feminist thoughts
  \item a women's college bulletin
  \item a collection of Shakespearean criticism
\end{choices}
\end{question}
\answer{3}
